% !TeX program = xelatex
% Faithful transcription of 23 handwritten pages; original images follow handwritten page order.
% Original medical statements and English spellings are retained.
\ifdefined\settextfont
  \def\NASStandalone{0}
\else
  \def\NASStandalone{1}
  \documentclass[11pt,a4paper]{article}
  \usepackage[margin=1.8cm]{geometry}
  \usepackage{amsmath,amssymb,enumitem}
  \usepackage{xepersian}
  \IfFontExistsTF{Amiri}{\settextfont{Amiri}}{\settextfont{FreeSerif}}
  \setlatintextfont{FreeSerif}
  \setlength{\parindent}{0pt}
\fi
\providecommand{\en}[1]{\lr{#1}}
\makeatletter
\newcommand{\NASalph}[1]{\@latinalph{\value{#1}}}
\makeatother
\newcommand{\NASarrow}{\en{$\longleftarrow$}}
\newcommand{\NASheader}{{\footnotesize\en{Note book\qquad Subject: 22\qquad Veor: \_\_\_\quad Month: \_\_\_\quad Date: \_\_\_}\par}}
\newcommand{\NASfooter}[1]{\par\medskip{\footnotesize\en{BAHAR DANESH ASL}\hfill #1\par}}
\ifnum\NASStandalone=1 \begin{document}\fi

% pic01.jpg = 20261019_083048149.jpg
\NASheader
\section*{\en{NonTraumatic abnormalities of the appendicular skeleton including arthritis}}
\begin{enumerate}[label=\arabic*.,leftmargin=*,itemsep=.5em]
\item اغلب بررسی‌ها استخوان با ۲ گرافی ساده عمود بر هم آغاز می‌شوند.
\item بهترین تصویربرداری برای بررسی مغز استخوان \NASarrow\ \en{MRI}
\item کورتکس استخوان در گرافی ساده بهتر از \en{MRI} دیده می‌شود.
\item \en{The shaft of the bone is the diaphysis, capped on each end by the epiphysis.}
\item محل اتصال دیافیز به اپی‌فیز \NASarrow\ متافیز؛ محل صفحهٔ رشد در کودکان.
\item در \en{MRI} زانو:
\begin{latin}\begin{itemize}[nosep]
\item Infrapatellar fat pad $\longrightarrow$ high signal (white)
\item Bone's cortex $\longrightarrow$ very weak signal (dark)
\end{itemize}\end{latin}
\item غضروف و مایع مفصلی در رادیوگرافی دیده نمی‌شوند.
\item \en{Bone abnormalities on conventional radiography}
\begin{latin}\begin{enumerate}[label=\alph*),nosep]
\item Pattern of increased bone density
\item Pattern of decreased bone density
\end{enumerate}\end{latin}
\end{enumerate}
\NASfooter{۱}

% pic02.jpg = 20261019_083053101.jpg
\clearpage\NASheader
\begin{enumerate}[label=\arabic*.,leftmargin=*,itemsep=.5em]\setcounter{enumi}{8}
\item \en{Bone density}
\begin{latin}\begin{itemize}
\item Increased
\begin{itemize}[nosep]
\item Diffuse: Diffuse osteoblastic metastases
\item Focal: Paget disease; Avascular necrosis of bone; Localized osteoblastic metastases
\end{itemize}
\item Decreased
\begin{itemize}[nosep]
\item Diffuse: Osteoporosis; Hyperparathyroidism
\item Focal: Localized osteolytic metastases; MM; Osteomyelitis
\end{itemize}
\end{itemize}\end{latin}
\item بیماری‌هایی که منجر به افزایش جنرالیزه دانسیته استخوان می‌شوند:
\begin{itemize}[nosep]\item مدولاری کویتی را \en{fill-in} می‌کنند و در رادیوگرافی هم دانسیتهٔ کورتکس دیده می‌شود.\end{itemize}
\item بیماری متاستاز استئوبلاستیک
\begin{itemize}[nosep]
\item کارسینوم پروستات \NASarrow\ افزایش جنرالیزه دانسیته استخوان
\item متاستازهای منتشر از سایر تومورها
\item همچنین می‌تواند افزایش دانسیته \en{focal} بدهند.
\item ضایعات استئوبلاستیک بیشتر به: مهره، دنده، لگن، هومروس، فمور
\end{itemize}
\end{enumerate}
\NASfooter{۲}

% pic03.jpg = 20261019_083100399.jpg
\clearpage\NASheader
\subsection*{\en{Focal sclerotic lesions}}
ضایعات اسکلروتیک ناحیه‌ای
\begin{itemize}[itemsep=.5em]
\item می‌توانند کورتکس، حفرهٔ مدولاری یا هر دو را درگیر کنند.
\item آن‌هایی که کورتکس را درگیر می‌کنند معمولاً:
\begin{latin}Periosteal new-bone formation $\longrightarrow$ Thickening of the cortex\end{latin}
\item آن‌هایی که حفرهٔ مدولاری را درگیر می‌کنند:
\begin{latin}Punctate, Amorphous, sclerotic lesions surrounded by the normal medullary cavity\end{latin}
\end{itemize}
\begin{enumerate}[label=\arabic*.,leftmargin=*,itemsep=.5em]\setcounter{enumi}{11}
\item تصویربرداری انتخابی برای شناسایی متاستازهای استخوانی
\begin{latin}Radionuclide bone scan\end{latin}
\item در متاستاز استخوانی منتشر در اسکن استخوان \NASarrow\ نمای \en{Superscan}
\item کارسینوم متاستاتیک برست \NASarrow\ می‌تواند هم استئولیتیک و هم استئوبلاستیک باشد.
\item متاستاز به استخوان از تومورهای اولیهٔ استخوان بسیار شایع‌تر است.
\item انواع متاستاز به استخوان: استئوبلاستیک، استئولیتیک؛ متاستازهای پستان هر دو ویژگی را دارند.
\item متاستاز به استخوان اگرچه هر بخشی می‌تواند باشد، اما اغلب دیستال آرنج یا زانو را درگیر نمی‌کند.
\begin{itemize}[nosep]\item اگر درگیر باشد اغلب متاستازها منتشر وجود دارد و منشأ معمولاً ریه یا برست است.\end{itemize}
\end{enumerate}
\NASfooter{۳}

% pic04.jpg = 20261019_083204512.jpg (Drive id 1IMyEp7Mjw_r9nys7yKm1NVR_HihpnOkC)
\clearpage\NASheader
\section*{نکروز آواسکولار استخوان (\en{AVN})}
\begin{enumerate}[label=\arabic*.,leftmargin=*,itemsep=.5em]
\item نام‌های دیگر: نکروز ایسکمیک، نکروز آسپتیک، استئونکروز
\item \en{AVN} معمولاً استخوان‌هایی را درگیر می‌کند که خون‌رسانی کولترال ضعیفی دارند.
\begin{itemize}[nosep]
\item اسکافوئید، سر فمور
\item در ابتدا اجزا خون‌ساز مغز استخوان از بین می‌روند.
\item \en{MRI} حساس‌ترین روش
\end{itemize}
\item در \en{AVN} در رادیوگرافی ساده بخش آواسکولار \en{Dense}تر می‌شود و نسبت به بقیهٔ استخوان اسکلروتیک دیده می‌شود.
\begin{itemize}[nosep]\item خصوصاً سر فمور و سر هومروس\end{itemize}
\item \en{AVN} سر فمور \NASarrow\ می‌تواند به مصرف طولانی‌مدت استروئید دیده شود.
\begin{itemize}[nosep]\item ساب‌کندرال لوسنسی ناشی از شکستگی‌های ساب‌کندرال \NASarrow\ \en{Crescent sign}\end{itemize}
\item \en{AVN} سر هومروس \NASarrow\ در بیماری سیکل سل؛ نمای \en{Snow-capping}
\end{enumerate}
\NASfooter{۴}

% pic05.jpg = 20261019_083112563.jpg
\clearpage\NASheader
\begin{enumerate}[label=\arabic*.,leftmargin=*,itemsep=.5em]\setcounter{enumi}{5}
\item در \en{AVN} در \en{MRI} عموماً یک کاهش سیگنال در محلی که قبلاً به علت مغز استخوان \en{fatty} وجود داشته دیده می‌شود.
\item \en{AVN in conventional radiography}
\begin{latin}Old medullary bone infarcts: Dense, amorphous deposits of bone within the medullary cavity of long bones. Frequently marginated by a thin, sclerotic membrane\end{latin}
\item علل \en{AVN}
\begin{itemize}[nosep]
\item انتراواسکولار: بیماری سایکل سل (\en{SCD})؛ پلی‌سایتمی ورا (\en{PV})
\item واسکولار: واسکولیت‌ها \NASarrow\ لوپوس؛ ناشی از رادیوتراپی
\item اکستراواسکولار: تروما
\item ایدیوپاتیک: مصرف استروئید، بیماری گوشه؛ \en{Legg-Calve-Perthes disease}
\end{itemize}
\item \en{MRI} حساس‌ترین روش بررسی \en{AVN} در هیپ است.
\item تشخیص افتراقی کلسیفیکیشن اینترامدولاری: انفارکت استخوان، \en{Enchondroma}
\end{enumerate}
\NASfooter{۵}

% pic06.jpg = 20261019_083118510.jpg
\clearpage\NASheader
\subsection*{\en{Finding metastases in Bone Scan}}
\begin{enumerate}[label=\arabic*.,leftmargin=*,itemsep=.5em]
\item در اسکن استخوان از تکنسیوم \en{99M} استفاده می‌شود.
\item ضایعات: استئوبلاستیک \NASarrow\ افزایش جذب؛ استئولیتیک \NASarrow\ اغلب افزایش جذب
\item اسکن استخوان در تشخیص \en{MM} حساسیت پایینی دارد.
\begin{itemize}[nosep]\item در شک به \en{MM} رادیوگرافی ساده درخواست می‌کنیم.\end{itemize}
\item اسکن استخوان ویژگی پایینی دارد.
\begin{itemize}[nosep]\item اگر مثبت شد \NASarrow\ رادیوگرافی، \en{CT}، \en{MRI}\end{itemize}
\end{enumerate}
\subsection*{بیماری پاژه}
\begin{enumerate}[label=\arabic*.,leftmargin=*,itemsep=.5em]
\item بیماری مزمن استخوان است.
\item بیشتر در مردان مسن
\item علت \NASarrow\ عفونت مزمن با پارامیکسوویروس
\item پاتولوژی \NASarrow\ درجات مختلف از افزایش بازجذب استخوان و افزایش همزمان ساخت استخوان
\item استخوان \en{Dense}تر می‌شود ولی ریسک شکستگی پاتولوژیک بالاتر می‌رود.
\end{enumerate}
\NASfooter{۶}

% pic07.jpg = 20261019_083126582.jpg
\clearpage\NASheader
\begin{enumerate}[label=\arabic*.,leftmargin=*,itemsep=.5em]\setcounter{enumi}{5}
\item بیشترین درگیری در لگن
\begin{itemize}[nosep]\item سپس \NASarrow\ لومبار اسپاین، توراسیک اسپاین، پروگزیمال فمور و کالواریوم\end{itemize}
\item تشخیص بیماری پاژه \NASarrow\ رادیوگرافی ساده
\begin{latin}\begin{itemize}[nosep]
\item Thickening of the cortex
\item Coarsening and thickening of the trabecular pattern
\item Increase in the size of the bone involved
\end{itemize}\end{latin}
\item یافتهٔ کلاسیک شرح حال \NASarrow\ کلاه دیگر اندازهٔ سر بیمار نیست.
\end{enumerate}
\section*{بیماری‌هایی که دانسیتهٔ استخوان را کاهش می‌دهند}
\begin{enumerate}[label=\arabic*.,leftmargin=*]
\item نمای کاهش دانسیتهٔ استخوان
\begin{latin}\begin{itemize}[itemsep=.4em]
\item Diffuse loss of the normal network of bony trabeculae in the medullary cavity.
\item Accentuation of the cortex
\begin{itemize}\item Cortex, although thinner than normal, stands out more strikingly because of the lower density of the medullary cavity.\end{itemize}
\item Compression of vertebral bodies
\item Pathologic fractures in the $\longrightarrow$ hip; pelvic; vertebral body
\end{itemize}\end{latin}
\end{enumerate}
\NASfooter{۷}

% pic08.jpg = 20261019_083131022.jpg
\clearpage\NASheader
\begin{enumerate}[label=\arabic*.,leftmargin=*,itemsep=.5em]\setcounter{enumi}{1}
\item استئوپروز
\begin{itemize}[itemsep=.4em]
\item انواع: پست‌منوپوزال؛ مرتبط با سن
\item پست‌منوپوزال استئوپروز: افزایش بازجذب استخوان به دلیل فعالیت استئوبلاستیک
\item استئوپروز مرتبط با سن: از سن ۴۵--۵۵ سالگی آغاز می‌شود.
\item عوامل خطری که ریسک استئوپروز را افزایش می‌دهند: استروئید اگزوژن، بیماری کوشینگ، کمبود استروژن، فعالیت فیزیکی ناکافی، الکلیسم
\item رادیوگرافی نرمال حساسیت کمی در تشخیص استئوپروز دارد؛ باید حداقل ۵۰\% دانسیتهٔ استخوان از دست رفته باشد.
\end{itemize}
\begin{latin}DEXA (Dual-energy X-ray Absorptiometry) Scan\end{latin}
\begin{itemize}[nosep]
\item دقیق‌ترین روش سنجش تراکم استخوان
\item ستون مهره و لگن بررسی می‌شود.
\item دوز اشعه بسیار کم
\end{itemize}
\end{enumerate}
\NASfooter{۸}

% pic09.jpg = 20261019_083140277.jpg
\clearpage\NASheader
\begin{enumerate}[label=\arabic*.,leftmargin=*,itemsep=.5em]\setcounter{enumi}{2}
\item هایپرپاراتیروئیدی
\begin{itemize}[nosep]
\item تولید بیش از حد \en{PTH} \NASarrow\ از دست رفتن تراکم استخوان
\item بهترین تصویربرداری پاراتیروئید: \en{US}، اسکن هسته‌ای، \en{MRI}
\end{itemize}
\item انواع هایپرپاراتیروئیدی
\begin{itemize}[itemsep=.4em]
\item هایپرپارای اولیه: عموماً به علت آدنوم پاراتیروئید؛ تقریباً همیشه هایپرکلسمی می‌دهد.
\item هایپرپارای ثانویه: هایپرپلازی پاراتیروئید به علت کمبود کلسیم و فسفر؛ عموماً در بیماری مزمن کلیه دیده می‌شود.
\item هایپرپارای ثالثیه: در پی هایپرپارای ثانویهٔ طولانی‌مدت؛ پاراتیروئید به‌طور خودمختار مقدار زیادی \en{PTH} آزاد می‌کند؛ همیشه همراه با هایپرکلسمی
\end{itemize}
\item یافته‌های هایپرپاراتیروئیدی در گرافی ساده
\begin{latin}\begin{itemize}[itemsep=.4em]
\item Overall decrease in bony density.
\item Subperiosteal bone resorption
\begin{itemize}\item on the radial side of the middle phalanges of the index and middle fingers.\end{itemize}
\item Erosion of the distal clavicles
\item Well-circumscribed lytic lesions in the long bones: Brown Tumors
\item Resorption of the terminal phalanges (acroosteolysis)
\end{itemize}\end{latin}
\end{enumerate}
\NASfooter{۹}

% pic10.jpg = 20261019_083144547.jpg
\clearpage\NASheader
\section*{کاهش ناحیه‌ای تراکم استخوان}
\begin{enumerate}[label=\arabic*.,leftmargin=*,itemsep=.5em]
\item متاستاز استئولیتیک
\begin{latin}Focal infiltration of bone by cells other than osteocytes\end{latin}
\begin{itemize}[itemsep=.4em]
\item مدولاری کویتی تقریباً همیشه درگیر می‌شود.
\item اگر صرفاً درگیری مدولاری کویتی رخ داده باشد باید حداقل ۵۰\% تودهٔ استخوانی از دست برود تا در گرافی ساده قابل تشخیص باشد؛ در نمای \en{en face}.
\item \en{MRI} حساس‌ترین روش تشخیص است.
\item در برخی بیماران صرفاً کورتکس درگیر می‌شود؛ در گرافی ساده راحت‌تر تشخیص داده می‌شود.
\end{itemize}
\item انواع الگوی ضایعات لیتیک استئولیز
\begin{latin}\begin{enumerate}[label=\alph*),itemsep=.5em]
\item Geographic bone lesion
\begin{itemize}\item gradual zone of transition\end{itemize}
\item Moth-eaten pattern (mottled)
\begin{itemize}[nosep]
\item ill-defined lytic lesions $\longrightarrow$ with indistinct margins
\item implying a more aggressive malignancy
\end{itemize}
\item Permeative pattern (Round cell lesions)
\begin{itemize}[nosep]
\item small irregular holes in the bone
\item E.x: Ewing's sarcoma, myeloma, leukemia
\end{itemize}
\end{enumerate}\end{latin}
\end{enumerate}
\NASfooter{۱۰}

% pic11.jpg = 20261019_083152043.jpg
\clearpage\NASheader
\begin{enumerate}[label=\arabic*.,leftmargin=*,itemsep=.5em]\setcounter{enumi}{2}
\item یافته‌ها متاستاز استئولیتیک در گرافی ساده
\begin{latin}\begin{itemize}[itemsep=.4em]
\item Irregular shaped, lucent bone lesions $\longrightarrow$ Single; Multiple
\item They can be expansile and soap-bubbly
\begin{itemize}\item Renal and Thyroid carcinoma\end{itemize}
\item In the spine $\longrightarrow$ Destroy the pedicles (because of their blood supply)
\begin{itemize}\item pedicle sign\end{itemize}
\end{itemize}\end{latin}
در افتراق متاستاز از \en{MM} کمک‌کننده است.
\item علل متاستازهای استئوبلاستیک
\begin{itemize}[nosep]
\item کارسینوم پروستات (شایع‌ترین علت در مردان مسن)
\item کارسینوم برست: عموماً استئولیتیک است ولی می‌تواند استئوبلاستیک نیز باشد، خصوصاً اگر درمان شود.
\item لنفوم
\item تومورهای کارسینوئید (نادر)
\end{itemize}
\item علل متاستازهای استئولیتیک
\begin{itemize}[nosep]
\item کنسر ریه (شایع‌ترین علت در مردان)
\item کنسر برست (شایع‌ترین علت در زنان)
\item \en{RCC}
\item کارسینوم تیروئید
\end{itemize}
\end{enumerate}
\NASfooter{۱۱}

% pic12.jpg = 20261019_083156392.jpg
\clearpage\NASheader
\section*{میلوما}
\begin{enumerate}[label=\arabic*.,leftmargin=*,itemsep=.5em]
\item شایع‌ترین بدخیمی اولیهٔ استخوان در بزرگسالان
\item \en{Can occur in a solitary form}
\begin{latin}\begin{itemize}[nosep]
\item often seen as a soap-bubbly, expansile lesion in $\longrightarrow$ Spine; Pelvis
\item Solitary Plasmacytoma
\end{itemize}\end{latin}
\item یافته‌ها \en{MM} در گرافی ساده
\begin{itemize}[itemsep=.4em]
\item شایع‌ترین تظاهر استئوپروز منتشر و عموماً شدید است.
\item پلاسماسیتوما
\begin{latin}\begin{itemize}[nosep]
\item Expansile
\item Septate lesion
\item Frequently associated with soft tissue masses
\end{itemize}\end{latin}
\item فرم منتشر
\begin{latin}\begin{itemize}[nosep]
\item Multiple
\item Small
\item sharply-circumscribed (punched-out) lytic lesions of the same size
\item without any sclerotic reaction around them
\end{itemize}\end{latin}
\end{itemize}
\end{enumerate}
\section*{استئومیلیت}
\begin{enumerate}[label=\arabic*.,leftmargin=*]
\item تخریب ناحیه‌ای استخوان \NASarrow\ عموماً به علت یک عامل عفونی منتقل از خون
\begin{itemize}[nosep]\item شایع‌ترین: استاف اورئوس\end{itemize}
\end{enumerate}
\NASfooter{۱۲}

% pic13.jpg = 20261019_083204512.jpg (Drive id 1Tegbq9CO2fq69FksrKvWqLgDTAGipbuq)
\clearpage\NASheader
\begin{enumerate}[label=\arabic*.,leftmargin=*,itemsep=.5em]\setcounter{enumi}{1}
\item استئومیلیت در اطفال بیشتر متافیز را درگیر می‌کند چون خون‌رسانی بیشتری دارد.
\item یافته‌ها استئومیلیت حاد در گرافی ساده
\begin{latin}\begin{itemize}[nosep]
\item Focal cortical bone destruction
\item Periosteal new bone formation
\item Soft tissue swelling and focal osteoporosis
\end{itemize}\end{latin}
\item در بزرگسالان عفونت تمایل دارد برخلاف کودکان، مفصل را نیز درگیر نماید که منجر به آرتریت سپتیک می‌شود.
\item رادیوگرافی ساده ۱۰ روز بعد از شروع استئومیلیت می‌تواند آن را تشخیص دهد.
\begin{itemize}[nosep]\item از \en{MRI} و اسکن هسته‌ای برای تشخیص زودهنگام استفاده می‌کنیم.\end{itemize}
\item در اسکن هسته‌ای، از اسکن \en{WBC} نشان‌دارشده با ایندیوم استفاده می‌شود.
\end{enumerate}
\section*{بیماری‌های مفصل}
\begin{enumerate}[label=\arabic*.,leftmargin=*,itemsep=.5em]
\item تصویربرداری اولیه برای بررسی آرتریت \NASarrow\ رادیوگرافی ساده
\item انواع آرتریت
\begin{itemize}[nosep]
\item آرتریت هایپرتروفیک
\item ایجاد استخوان در محل مفصل درگیر
\item انواع: \en{Subchondral sclerosis}؛ \en{Osteophyte}
\end{itemize}
\end{enumerate}
\NASfooter{۱۳}

% pic14.jpg = 20261019_083210237.jpg
\clearpage\NASheader
\begin{itemize}[itemsep=.5em]
\item آرتریت ایروزیو \NASarrow\ بیانگر التهاب زمینه‌ای
\begin{latin}Tiny, marginal, irregularly shaped lytic lesions in or around the joint surface. $\longrightarrow$ Called Erosions\end{latin}
\item آرتریت عفونی \NASarrow\ تورم مفصل، استئوپنی
\begin{latin}Destruction of long, contiguous segments of the articular cortex.\end{latin}
\end{itemize}
\begin{enumerate}[label=\arabic*.,leftmargin=*,itemsep=.5em]\setcounter{enumi}{2}
\item انواع آرتریت هایپرتروفیک
\begin{enumerate}[label=\en{\NASalph{enumii})},itemsep=.4em]
\item استئوآرتریت اولیه: شایع‌ترین علت، استرس مکانیکی؛ دست‌ها، هیپ، زانو
\item استئوآرتریت ثانویه: \en{DJD} ثانویه به تروما؛ \en{AVN}
\item آرتروپاتی شارکو: عموماً ثانویه به \en{DM}؛ فرگمنتیشن، تخریب مفصل و اسکلروز
\item \en{CPPD}: کندروکلسینوز
\end{enumerate}
\item انواع آرتریت عفونی
\begin{enumerate}[label=\en{\NASalph{enumii})},nosep]
\item توبرکلوز، پیوژنیک \NASarrow\ همراه با استئوپنی
\end{enumerate}
\end{enumerate}
\NASfooter{۱۴}

% pic15.jpg = 20261019_083218293.jpg
\clearpage\NASheader
\begin{enumerate}[label=\arabic*.,leftmargin=*]\setcounter{enumi}{4}
\item انواع آرتریت ایروزیو
\begin{enumerate}[label=\en{\NASalph{enumii})},itemsep=.6em]
\item روماتوئید \NASarrow\ \en{MCP}، \en{PIP}، استئوپروز، تورم
\item \en{Gout} \NASarrow\ ایروژن جاکستا آرتیکولار؛ \en{MTP} انگشت اول؛ بدون استئوپروز
\item پسوریاتیک \NASarrow\ ایروژن جاکستا آرتیکولار مفصل اینترفالنژیال؛ دفورمیتی \en{Pencil in cup}
\item هموفیلیا \NASarrow\ ریمودلینگ در پی همارتروز؛ تغییرات مشابه در زانو بیماران فاقد هم:
\begin{latin}Think of juvenile rheumatoid arthritis\end{latin}
\item اسپوندیلیت آنکیلوزان
\begin{itemize}[nosep]
\item \en{HLA Ag B27 positive}
\item ساکروایلیت دوطرفه
\item ستون فقرات
\end{itemize}
\item اسپوندیلوآرتروپاتی سرونگاتیو
\begin{itemize}[nosep]
\item \en{RF Negative}
\item \en{HLA Ag B27 positive}
\item ستون فقرات
\item ساکروایلیت دوطرفه
\item آرتریت ری‌اکتیو
\item پسوریازیس
\end{itemize}
\end{enumerate}
\end{enumerate}
\NASfooter{۱۵}

% pic16.jpg = 20261019_083222389.jpg
\clearpage\NASheader
\begin{enumerate}[label=\arabic*.,leftmargin=*,itemsep=.5em]\setcounter{enumi}{5}
\item یافته‌ها تصویربرداری استئوآرتریت اولیه
\begin{latin}\begin{itemize}[itemsep=.3em]
\item Marginal osteophyte formation
\item Subchondral sclerosis
\item Subchondral cyst
\begin{itemize}[nosep]
\item Result of chronic impaction and necrosis
\item Varying in size
\end{itemize}
\item Narrowing of the joint space
\begin{itemize}\item Seen in all forms of arthritis.\end{itemize}
\end{itemize}\end{latin}
\item در اکثر بیماران استئوآرتریت مفاصل اینترفالنژیال دست
\begin{itemize}[nosep]\item مفصل کارپومتاکارپال اول نیز درگیر است.\end{itemize}
\item استئوآرتریت به‌طور شایع مفاصل دیستال اینترفالنژیال را درگیر می‌کند، خصوصاً در خانم‌های مسن.
\item یافته‌ها تصویربرداری استئوآرتریت ثانویه مشابه اولیه است با این تفاوت که:
\begin{itemize}[nosep]
\item در سنین غیرمعمول رخ می‌دهد.
\item ظاهر غیرمعمول دارد؛ عموماً یک‌طرفه است.
\item ممکن است در مکان غیرمعمول باشد؛ مثلاً آرنج.
\end{itemize}
\item آرتریت ایروزیو بیشتر در خانم‌های پره‌منوپوز رخ می‌دهد.
\item نمای \en{Gull-wing deformity} در مفاصل \en{PIP} در آرتریت ایروزیو
\end{enumerate}
\NASfooter{۱۶}

% pic17.jpg = 20261019_083229025.jpg
\clearpage\NASheader
\begin{enumerate}[label=\arabic*.,leftmargin=*,itemsep=.5em]\setcounter{enumi}{11}
\item مفاصل شایع برای آرتریت ایروزیو: \en{DIP, PIP}؛ کارپومتاکارپال اول؛ \en{IP} شست
\item در آرتریت ایروزیو ممکن است آنکیلوز استخوانی رخ دهد.
\item آرتروپاتی شارکو (مفصل نوروپاتیک)
\begin{itemize}[nosep]
\item اختلال در حس مفصل \NASarrow\ بیمار درد را کمتر احساس می‌کند.
\item شکستگی‌های کوچک متعدد
\item تظاهر اصلی بیماری التهاب و تورم مفصل است.
\item شایع‌ترین مورد در بیماران \en{DM} است.
\item بیشتر اندام تحتانی را درگیر می‌کند \NASarrow\ مچ پا و پا
\end{itemize}
\item علل آرتروپاتی شارکو
\begin{itemize}[nosep]
\item شانه: سیفلیس؛ تومور نخاعی؛ \en{Syrinx}
\item هیپ: \en{DM}؛ سیفلیس ثالثیه
\item مچ پا و پا: \en{DM} (شایع)؛ سیفلیس (ناشایع)
\end{itemize}
\item آرتروپاتی شارکو در \en{DDX} استئومیلیت قرار دارد چون هر دو تخریب استخوانی و واکنش پری‌اُست دارند.
\begin{itemize}[nosep]\item افتراق با \en{WBC Scan}\end{itemize}
\end{enumerate}
\NASfooter{۱۷}

% pic18.jpg = 20261019_083234236.jpg
\clearpage\NASheader
\begin{enumerate}[label=\arabic*.,leftmargin=*,itemsep=.5em]\setcounter{enumi}{16}
\item یافته‌ها رادیوگرافی آرتروپاتی شارکو
\begin{latin}\begin{itemize}[itemsep=.4em]
\item Extensive subchondral sclerosis
\begin{itemize}\item As a hypertrophic arthritis\end{itemize}
\item Fragmentation of bones surrounding the joint
\begin{itemize}\item May be resorbed and no longer visible\end{itemize}
\item Destruction of the joint
\end{itemize}\end{latin}
\item هالمارک دفورمیتی در شارکو فوت
\begin{latin}Collapse of the midfoot, called a rocket-bottom foot\end{latin}
\end{enumerate}
\section*{بیماری کلسیم پیروفسفات دپوزیشن (\en{CPPD})}
\begin{enumerate}[label=\arabic*.,leftmargin=*,itemsep=.5em]
\item ناشی از رسوب کریستال‌های \en{CPPD} در داخل و اطراف مفصل
\item عمدتاً در غضروف هیالینی، فیبروکارتیلج
\item در غضروف منیسک فیبروکارتیلج مچ دست و منیسک‌های زانو شایع است.
\item اسامی متفاوت برای آن
\begin{enumerate}[label=\en{\NASalph{enumii})}]
\item کندروکلسینوز \NASarrow\ کلسیفیه شدن غضروف مفصلی
\begin{itemize}[nosep]
\item در ۵۰\% افراد بالای ۸۵ سال دیده می‌شود که عموماً بی‌علامت است.
\item در بیماری‌های دیگر نیز دیده می‌شود: هایپرپارا، هموکروماتوز
\end{itemize}
\end{enumerate}
\end{enumerate}
\NASfooter{۱۸}

% pic19.jpg = 20261019_083241539.jpg
\clearpage\NASheader
\begin{enumerate}[label=\en{\NASalph{enumi})},leftmargin=*,itemsep=.6em]\setcounter{enumi}{1}
\item سودوگوت \NASarrow\ یک سندرم بالینی است.
\begin{latin}Acute monoarticular arthropathy\end{latin}
\begin{itemize}[nosep]
\item درد، قرمزی، تورم
\item شایع‌ترین مفصل: زانو
\item تشخیص با آسپیراسیون
\end{itemize}
\item آرتروپاتی پیروفسفات \NASarrow\ یک تشخیص رادیولوژیک است؛ شایع‌ترین فرم \en{CPPD} است.
\end{enumerate}
\begin{enumerate}[label=\arabic*.,leftmargin=*]\setcounter{enumi}{4}
\item آرتروپاتی پیروفسفات ممکن است از استئوآرتریت ثانویه غیرقابل افتراق باشد. به تفاوت‌ها توجه کنید:
\begin{itemize}[itemsep=.4em]
\item \en{CPPD} در مفاصل: پاتلوفمورال، رادیوکارپال، متاکارپوفالنژیال، مچ دست
\item در \en{CPPD} عموماً کندروکلسینوز وجود دارد.
\item کیست‌های ساب‌کندرال در \en{CPPD}: شایع‌تر؛ بزرگ‌تر؛ تعداد بیشتر
\end{itemize}
\begin{latin}\begin{itemize}[itemsep=.4em]
\item Hook-shaped bony excrescences along with second and third metacarpal heads in CPPD
\item Separation of the scaphoid and the lunate by more than 3 mm $\longrightarrow$ SLD
\item Collapse of the distal carpal row toward the radius $\longrightarrow$ SLAC
\end{itemize}\end{latin}
\end{enumerate}
\NASfooter{۱۹}

% pic20.jpg = 20261019_083247449.jpg
\clearpage\NASheader
\section*{آرتریت ایروزیو}
\begin{enumerate}[label=\arabic*.,leftmargin=*,itemsep=.5em]
\item آرتریت روماتوئید (\en{RA}) \NASarrow\ در زنان شایع‌تر است.
\begin{itemize}[nosep]
\item به‌طور شایع مفاصل پروگزیمال دست و مفصل مچ را درگیر می‌کند.
\item عموماً دوطرفه و قرینه است.
\item گرافی ساده تصویربرداری اولیه برای بررسی \en{RA} می‌باشد.
\item اولین یافته تورم بافت نرم و استئوپروز است.
\item استئوپروز در دو طرف مفصل شدیدتر است \NASarrow\ \en{Periarticular osteoporosis}
\end{itemize}
\item \en{RA} در دست
\begin{itemize}[nosep]
\item بیشتر مفاصل پروگزیمال را درگیر می‌کند: \en{CMC}، \en{MCP}، \en{PIP}
\item یافته‌های تأخیری
\end{itemize}
\begin{latin}\begin{itemize}[nosep]
\item Ulnar deviation of the fingers at MCP
\item Subluxation of the MCP and ligamentous laxity
\begin{itemize}[nosep]\item Swan-neck; Boutonnière\end{itemize}
\end{itemize}\end{latin}
\item \en{RA} در مچ دست
\begin{latin}\begin{itemize}[itemsep=.4em]
\item Erosions of the carpal, ulnar styloid and narrowing of the radiocarpal joint
\item Larger joints usually show no erosions but there may be narrowing of the joint space with no subchondral sclerosis
\end{itemize}\end{latin}
\end{enumerate}
\NASfooter{۲۰}

% pic21.jpg = 20261019_083305936.jpg
\clearpage\NASheader
\begin{enumerate}[label=\arabic*.,leftmargin=*,itemsep=.5em]\setcounter{enumi}{3}
\item \en{RA} در ستون فقره‌ها
\begin{itemize}[nosep]\item عمدتاً بخش سرویکال را درگیر می‌کند.\end{itemize}
\begin{latin}\begin{itemize}[itemsep=.4em]
\item Ligamentous laxity $\longrightarrow$ Forward subluxation of C$_1$ and C$_2$ (atlantoaxial subluxation)
\begin{itemize}\item Can produce Cord compression\end{itemize}
\item May also cause narrowing, sclerosis and fusion of the facet joints.
\end{itemize}\end{latin}
\item \en{Gout}
\begin{itemize}[nosep]
\item تغییرات التهابی مفصل به دنبال تجمع کریستال‌های کلسیم اورات
\item بین شروع علائم و مشاهدهٔ تغییرات استخوانی ۵--۷ سال فاصله است.
\item تشخیص \en{Gout} بالینی است نه رادیولوژیک.
\item عموماً در ابتدا مونوآرتیکولار است.
\item درگیری غیرقرینه
\item در مردان شایع‌تر است.
\end{itemize}
\item یافته‌ها تصویربرداری \en{Gout}
\begin{latin}\begin{itemize}[itemsep=.4em]
\item Halmark: Sharply marginated juxtaarticular erosion that tends to have a sclerotic border. (Rat-bites)
\item Joint space narrowing (Late finding)
\item Little or no periarticular osteoporosis
\item Tophi (Late finding, rarely calcify)
\item Olecranon bursitis (common)
\end{itemize}\end{latin}
\end{enumerate}
\NASfooter{۲۱}

% pic22.jpg = 20261019_083311523.jpg
\clearpage\NASheader
\begin{enumerate}[label=\arabic*.,leftmargin=*,itemsep=.5em]\setcounter{enumi}{6}
\item آرتریت پسوریاتیک
\begin{itemize}[nosep]
\item اغلب بیماران پسوریازیس سال‌ها قبل از آرتریت دچار تغییرات پوست و ناخنی می‌شوند.
\item بیماران پسوریازیس بلندمدت ۲۵\% احتمال دارد دچار آرتریت شوند.
\item معمولاً درگیری به‌صورت الیگوآرتیکولار است.
\item مفاصل درگیر: مفاصل کوچک دست، خصوصاً \en{DIP}
\end{itemize}
\item یافته‌های تصویربرداری آرتریت پسوریاتیک
\begin{latin}\begin{itemize}[nosep]
\item Juxtaarticular erosions (especially DIP)
\item Bony proliferation at the sites of tendon insertions (Enthesophytes)
\item Periosteal reaction (not common)
\item Pencil-in-cup deformity
\item Absence of osteoporosis
\end{itemize}\end{latin}
\item آرتریت پسوریاتیک در مفصل ساکروایلیاک
\begin{itemize}[nosep]
\item دوطرفه ولی غیرقرینه ساکروایلیت می‌دهد.
\item برخلاف اسپوندیلیت آنکیلوزان مفصل کاملاً \en{fuse} نمی‌شود.
\end{itemize}
\end{enumerate}
\section*{آرتریت سپتیک}
بهبود آرتریت سپتیک با فیبروز و آنکیلوز استخوانی رخ می‌دهد.
\begin{enumerate}[label=\arabic*.,leftmargin=*,itemsep=.5em]
\item انتقال عامل عفونی هماتوژن است.
\item انواع آرتریت سپتیک
\begin{itemize}[nosep]
\item پیوژنیک: استافیلوکوک، گنوکوک
\item نان‌پیوژنیک: \en{TB}
\end{itemize}
\end{enumerate}
\NASfooter{۲۲}

% pic23.jpg = 20261019_083318380.jpg
\clearpage\NASheader
\begin{enumerate}[label=\arabic*.,leftmargin=*,itemsep=.5em]\setcounter{enumi}{2}
\item ریسک‌فاکتورها آرتریت سپتیک
\begin{itemize}[nosep]
\item پروتز مفصلی، جراحی، تروما
\item \en{IV Drug use}
\item استروئید
\end{itemize}
\item محل‌های شایع آرتریت سپتیک
\begin{itemize}[nosep]
\item کودکان: زانو، هیپ
\item بزرگسالان: زانو
\item دست: براساس گازگرفتن یا دیابت
\end{itemize}
\item اگرچه گرافی ساده معمولاً اولین تصویربرداری است ولی حساسیت کمی دارد.
\begin{itemize}[nosep]\item یافته‌ها اولیه شامل تورم و استئوپنی است.\end{itemize}
\item یافته‌ها تصویربرداری آرتریت سپتیک
\begin{latin}\begin{itemize}[nosep]
\item Distruction of the articular cartilage
\begin{itemize}\item Distraction is usually relatively rapid\end{itemize}
\end{itemize}\end{latin}
\item آرتریت سپتیک عموماً مونوآرتیکولار است.
\item اگر قویاً به آرتریت سپتیک شک داریم \NASarrow\ آسپیراسیون
\item \en{MRI} برای تعیین آرتریت سپتیک بسیار حساس است.
\begin{latin}\begin{itemize}[nosep]
\item Enhancement of the synovium
\item Presence of a joint effusion
\end{itemize}\end{latin}
\item نوع نان‌پیوژنیک
\begin{itemize}[nosep]
\item \en{gradual loss of joint space}
\item عموماً مونوآرتیکولار
\item استئوپروز شدید شایع است.
\item در کودکان \en{Spine} و در بزرگسالان زانو
\end{itemize}
\end{enumerate}
\NASfooter{۲۳}

\ifnum\NASStandalone=1 \end{document}\fi
